\scanpage{120}
% archive foliation: 119
\uline{Invariant} sous l'action de $f^{-1}(H_S)$ i.e. telle que
\[
G(g\cdot s,g'\cdot t)=G(s,t)\quad\text{si}\quad
\left\{\begin{array}{l}s,t\in\Gamma(U,\mathcal P_X)\\g,g'\in\Gamma(V,H_V)\end{array}\right.
\]
\diagram{\begin{tikzcd}
U\arrow[r,"\text{ét}"]\arrow[d,no head]&X\arrow[d,"f"]\\
V\arrow[r,no head,"\text{ét}"']&S
\arrow[ul,phantom,"\text{comm.}"]
\end{tikzcd}}
et satisfaisant
\[
G(s,t)+G(t,r)=G(s,r).
\]
Cela permet donc de définir le faisceau \inserted{$\operatorname{Conn}_{X/S}$} (ét) des conn. projectives sur $X$ rel. à $S$, qui est un torseur sous $\underline\omega_{X/S}^{\otimes2}$, et d'interpréter $G(t,s)$ à l'aide d'un morphisme canonique
\[
c^X:\mathcal P_{X/S}\longrightarrow\operatorname{Conn}_{X/S}
\]
par
\[
G(t,s)=c^X_{\unsure{t}}(s)-c^X(t).
\]

Enfin, par construction la formation de $\mathcal P_{X/S}$, $\operatorname{Conn}_{X/S}$ et $c^X$ commute à la loc. étale. Cela permet d'interpréter encore
\[
c^X(t)=t^*(c)
\]
où $c$ est l'unique connexion sur $\mathbb P^1_S$.

Plus gén., pour un fibré en droites projectives $X_0$ sur $S$, il y a une unique conn. proj. $c_{X_0}$ sur $X_0$, et si $t,s:X\longrightarrow X_0$ sont étales on peut définir
\[
G(t,s)=c_{(s)}-c_{(t)}=s^*(c_{X_0})-t^*(c_{X_0})
\]
satisfaisant aux conditions analogues à celles du cas $X_0=\mathbb P^1_S$, que le lecteur (bis) \unsure{énoncera}.
